\documentclass[11pt,a4paper]{article}

%====================================================
% PACKAGES
%====================================================
\usepackage[utf8]{inputenc}
\usepackage[T1]{fontenc}
\usepackage[french]{babel}
\usepackage[a4paper,margin=2cm]{geometry}
\usepackage{amsmath,amssymb,amsthm,mathtools}
\usepackage{lmodern}

%====================================================
% ENVIRONNEMENTS
%====================================================
\newtheorem{definition}{Définition}
\newtheorem{proposition}{Proposition}
\newtheorem{theorem}{Théorème}
\newtheorem{exemple}{Exemple}
\newtheorem{remarque}{Remarque}

%====================================================
\begin{document}

\title{\textbf{Chapitre — Probabilités et Statistiques}}
\maketitle

\tableofcontents

%====================================================
\section{Pourquoi ce chapitre est-il fondamental ?}

Les probabilités permettent de modéliser l’incertitude.

\textbf{Idée :} transformer le hasard en modèle mathématique.

\textbf{Point clé :} une probabilité est une modélisation.

%====================================================
\section{Modélisation probabiliste}

\subsection{Univers}

\begin{definition}
Un univers $\Omega$ est l’ensemble des résultats possibles.
\end{definition}

\begin{exemple}
\[
\Omega = \{1,2,3,4,5,6\}
\]
\end{exemple}

\subsection{Événements}

\begin{definition}
Un événement est une partie de $\Omega$.
\end{definition}

%====================================================
\section{Tribu et modèle}

\begin{definition}
Une tribu $\mathcal{T}$ est une famille stable par complémentaire et union.
\end{definition}

%====================================================
\section{Probabilité}

\begin{definition}
\[
P : \mathcal{T} \to [0,1]
\]
avec $P(\Omega)=1$ et additivité.
\end{definition}

\begin{definition}
Un espace probabilisé est $(\Omega,\mathcal{T},P)$.
\end{definition}

%====================================================
\section{Propriétés}

\begin{proposition}
\[
P(A^c)=1-P(A)
\]
\end{proposition}

\begin{proposition}
\[
P(A \cup B)=P(A)+P(B)-P(A \cap B)
\]
\end{proposition}

%====================================================
\section{Probabilité conditionnelle}

\begin{definition}
\[
P(A|B)=\frac{P(A \cap B)}{P(B)}
\]
\end{definition}

%====================================================
\section{Indépendance}

\begin{definition}
\[
P(A \cap B)=P(A)P(B)
\]
\end{definition}

%====================================================
\section{Variables aléatoires}

\begin{definition}
\[
X : \Omega \to \mathbb{R}
\]
\end{definition}

%====================================================
\section{Espérance et variance}

\begin{definition}
\[
E(X)=\sum xP(X=x)
\]
\end{definition}

\begin{definition}
\[
Var(X)=E(X^2)-E(X)^2
\]
\end{definition}

%====================================================
\section{Lois discrètes usuelles}

\subsection{Bernoulli}

\[
P(X=1)=p,\quad P(X=0)=1-p
\]

\subsection{Binomiale}

\[
P(X=k)=\binom{n}{k}p^k(1-p)^{n-k}
\]

\subsection{Géométrique}

\[
P(X=k)=(1-p)^{k-1}p
\]

\subsection{Poisson}

\[
P(X=k)=\frac{\lambda^k e^{-\lambda}}{k!}
\]

%====================================================
\section{Lien entre les lois}

\subsection{Schéma de Bernoulli}

Suite d’expériences indépendantes de probabilité $p$.

\subsection{Bernoulli → Binomiale}

\begin{theorem}
Nombre de succès :
\[
X \sim \mathcal{B}(n,p)
\]
\end{theorem}

\begin{proof}
Choix des succès :
\[
\binom{n}{k}
\]
Probabilité :
\[
p^k(1-p)^{n-k}
\]
\end{proof}

\subsection{Binomiale → Poisson}

\begin{theorem}
Si $np \to \lambda$ :
\[
\mathcal{B}(n,p) \to \mathcal{P}(\lambda)
\]
\end{theorem}

%====================================================
\section{Espérance et variance des lois}

\subsection{Bernoulli}

\[
E(X)=p,\quad Var(X)=p(1-p)
\]

\begin{proof}
Calcul direct.
\end{proof}

\subsection{Binomiale}

\[
E(X)=np,\quad Var(X)=np(1-p)
\]

\begin{proof}
Somme de variables Bernoulli indépendantes.
\end{proof}

\subsection{Géométrique}

\[
E(X)=\frac{1}{p},\quad Var(X)=\frac{1-p}{p^2}
\]

\subsection{Poisson}

\[
E(X)=\lambda,\quad Var(X)=\lambda
\]

%====================================================
\section{Variables continues}

\begin{definition}
\[
P(a \leq X \leq b)=\int_a^b f(x)dx
\]
\end{definition}

%====================================================
\section{Lois continues}

\subsection{Uniforme}

\[
f(x)=\frac{1}{b-a}
\]

\subsection{Exponentielle}

\[
f(x)=\lambda e^{-\lambda x}
\]

\subsection{Normale}

\[
Z=\frac{X-\mu}{\sigma}
\]

%====================================================
\section{Statistiques}

\subsection{Échantillon}

Variables indépendantes.

\subsection{Intervalle de confiance}

\[
\bar{x} \pm \frac{1}{\sqrt{n}}
\]

\subsection{Fluctuation}

Comparer fréquence et probabilité.

%====================================================
\section{Résumé}

\begin{itemize}
\item modèle $(\Omega,\mathcal{T},P)$
\item lois usuelles
\item espérance et variance
\item statistiques
\end{itemize}

\end{document}